\section{Unfolding the First Power Allocation}
The proposal network acts upstream of an exact allocation. A different opportunity arises inside the allocation itself: an inexpensive direction can be discarded by the LP, whereas a better power iterate can immediately improve the design under a tight iteration budget. We therefore test a second learning mechanism with the fixed transmit proposal and the same analytic receiver.

\subsection{An eight-layer primal--dual allocator}
Use $x$ for the power fractions of all atoms, $a=\bar\alpha$ for their scaled rewards and $\phi(x)=a^\top x/(1+c^\top x)$ for the scaled objective. Here $a$ is real and stacks rewards across all tones, while $a_n$ in the signal model is a complex steering vector on one tone. Equilibrate each user row of $B$ and its entry in $d$ by the maximum absolute coefficient in that row. Let $z\ge0$ collect one multiplier for each of the $NK$ user constraints. Initialize the best feasible allocation $x_b$ with zero-forcing user powers and distribute the surplus equally over the null-space sensing directions. Set $x^0=x_b$, $z^0=0$ and $r=\phi(x_b)$. For layers $l=0,\ldots,7$, compute
\begin{align}
 x^{l+1}&=\Pi_{\Delta}\!\left(x^l+\tau_l(a-r c+B^\top z^l)\right),\label{eq:pd1}\\
 z^{l+1}&=\left[z^l+\upsilon_l\{d-B(2x^{l+1}-x^l)\}\right]_+ .
 \label{eq:pd2}
\end{align}
Here $[\cdot]_+$ takes the componentwise positive part, $\Delta=\{x\ge0:\mathbf1^\top x=1\}$ and $\Pi_\Delta$ is Euclidean projection. The positive scalars $\tau_l$ and $\upsilon_l$ control the changes in power and user prices, respectively. The gradient in \eqref{eq:pd1} follows by differentiating the affine fractional-programming residual $a^\top x-r(1+c^\top x)$ plus $z^\top(Bx-d)$. Equation~\eqref{eq:pd2} raises the price of an extrapolated user deficit. A power constraint of equality is appropriate for these raw iterates: scaling a feasible, positive-reward allocation upwards to exhaust its budget preserves $Bx\ge d$ and increases $\phi$ because the denominator contains the constant one.

After each layer, decode a feasible copy of $x^{l+1}$ as in Appendix~\ref{ap:decode}. If that copy has a higher design score, replace $x_b$ and $r$ by the copy and its score. The next raw update starts from $x^{l+1}$. Return $x_b$ and the last price vector. Physical-stream recovery and its audit are shared with the exact LP. Unlike an LP multiplier, a finite-depth price is approximate; when used to enrich a bank it supplies a proposed direction. In the six-allocation hybrid, later allocations use exact LP duals.

\subsection{What is trained}
The tied control uses the same pair $(\tau_l,\upsilon_l)=(0.5,5.6)$ in every layer. The learned allocator uses four successive parameter pairs, each repeated for two layers. Their trained values in order are $(0.125,7.84)$, $(0.0625,5.6)$, $(4,1.4)$ and $(5.6,22.4)$. These eight real scalars are shared across channels, loads and tones. We choose them by multiplicative coordinate search to maximize the mean returned design SCNR in dB over the training cases. The tied control is itself tuned by grid search and coordinate refinement. Validation chooses between the initial tied schedule and three passes of untied training. The channel seeds are 31001--31008 for training, 32001--32004 for validation, and 33001--33012 for testing. Training and validation use $K\in\{4,8,12\}$ and SINR requirements of 4 and 16 dB. This experiment learns solver steps from estimated-kernel objectives; it uses the ideal kernel only for evaluation.
